\documentclass[11pt]{article}
\usepackage[T1]{fontenc}
\usepackage{amsmath,amssymb,amsthm,lmodern}
\usepackage[margin=1in]{geometry}
\usepackage{hyperref}
\hypersetup{hidelinks,hypertexnames=false,pdftitle={MCMT: Optimal Strong Stationary Times, Proposition 6.21},pdfauthor={David A. Levin; Yuval Peres}}
\newcommand{\X}{\mathcal{X}}
\newcommand{\Ft}{\mathcal{F}_t}
\newcommand{\Px}{\mathbf{P}_x}
\newcommand{\sourceloc}[1]{\par\smallskip\noindent\textit{Source: #1.}\par\nopagebreak[4]\smallskip}
\newcommand{\editorial}[1]{\par\smallskip\noindent\textsf{Editorial note.} #1\par\smallskip}
\renewcommand{\qedsymbol}{\ensuremath{\blacksquare}}
\title{Optimal Strong Stationary Times\\Proposition 6.21}
\author{David A. Levin \and Yuval Peres}
\date{Private reading transcription of the fixed second-edition author PDF}
\begin{document}
\maketitle

% BEGIN CONTEXT MCMT:6.21
\noindent\textbf{Source and scope.}
\emph{Markov Chains and Mixing Times, second edition}, by David A. Levin
and Yuval Peres, ``With contributions by Elizabeth L. Wilmer'' as printed
on the cover, physical page 1. The fixed author-hosted PDF has 461 physical
pages and PDF generation metadata dated March 21, 2017; a separate author
revision date is not stated on the inspected cover.
The source is \href{https://pages.uoregon.edu/dlevin/MARKOV/mcmt2e.pdf}{\nolinkurl{https://pages.uoregon.edu/dlevin/MARKOV/mcmt2e.pdf}}.
The sole primary selection is Section 6.7, Proposition 6.21 and its complete
printed proof on physical page 101, printed page 85. The selection is in
the book's finite-chain setting. Selected definitions, identities and the
original Exercise 6.4 below are uncounted context. This transcription
remains private.

\sourceloc{Selected original Notes, physical page 103, printed page 87}
Strong stationary times were introduced by Aldous and Diaconis (1986,
1987). An important class of strong stationary times was constructed by
Diaconis and Fill (1990). The thesis of Pak (1997) has many examples of
strong stationary times.

Section 6.7 comes from Aldous and Diaconis (1987).

The strong stationary time we give for the cycle in Section 6.5.5 is due
to Diaconis and Fill (1990), although the exposition is different. The idea
goes back to the construction of a Skorokhod embedding due to Dubins (1968).

\editorial{These are the book's attribution statements. The printed
Section 6.7 proof is reproduced as the book's exposition; the cycle
construction mentioned in the Notes is outside this selection. The
original time ranges and quotient formulas below are retained, with no
additional endpoint or zero-denominator conventions.}

\subsection*{Selected definitions and identities}
\sourceloc{Section 6.2, physical page 93, printed page 77}
Let $\{\Ft\}$ be a \emph{filtration}, a sequence of $\sigma$-algebras such
that $\Ft\subset\mathcal{F}_{t+1}$ for all $t$. We say that $\{X_t\}$ is
\emph{adapted} to $\{\Ft\}$ if $X_t$ is $\Ft$-measurable for all $t$.
If $\mathcal{H}_t=\sigma(X_0,X_1,\ldots,X_t)$, then
$\{\mathcal{H}_t\}$ is called the \emph{natural filtration}. Clearly
$\{X_t\}$ is adapted to the natural filtration.

Suppose that $\{X_t\}$ is adapted to $\{\Ft\}$. We say that $\{X_t\}$ is
a \emph{Markov chain with respect to $\{\Ft\}$} if
\begin{equation}
\Px\{X_{t+1}=y\mid\Ft\}=P(X_t,y),
\tag{6.1}\label{eq:6.1}
\end{equation}
where $P$ is a transition matrix. A Markov chain as defined by (1.1)
satisfies \eqref{eq:6.1} when $\{\Ft\}$ is the natural filtration.

A \emph{stopping time} for the filtration $\{\Ft\}$ is a random variable
$\tau$ with values in $\{0,1,\ldots\}$ such that $\{\tau=t\}\in\Ft$.
If the filtration of a stopping time is not specified, it will be assumed
to be the natural filtration.

\sourceloc{Section 6.3, physical page 94, printed page 78; definition only}
Let $(X_t)$ be an irreducible Markov chain with stationary distribution
$\pi$. Suppose that $\{\Ft\}$ is a filtration, and $\{X_t\}$ is adapted
to $\{\Ft\}$. A \emph{stationary time} $\tau$ for $(X_t)$ is an
$\{\Ft\}$-stopping time, possibly depending on the starting position $x$,
such that the distribution of $X_\tau$ is $\pi$:
\begin{equation}
\Px\{X_\tau=y\}=\pi(y)\qquad\text{for all }y.
\tag{6.2}\label{eq:6.2}
\end{equation}

\sourceloc{Section 6.4 and full Remark 6.8, physical page 94, printed page 78}
Let $(X_t)$ be a Markov chain with respect to the filtration $\{\Ft\}$,
with stationary distribution $\pi$. A \emph{strong stationary time} for
$(X_t)$ and starting position $x$ is an $\{\Ft\}$-stopping time $\tau$,
such that for all times $t$ and all $y$,
\begin{equation}
\Px\{\tau=t,X_\tau=y\}=\Px\{\tau=t\}\pi(y).
\tag{6.3}\label{eq:6.3}
\end{equation}
In words, $X_\tau$ has distribution $\pi$ and is independent of $\tau$.

\noindent\textbf{Remark 6.8.} If $\tau$ is a strong stationary time starting
from $x$, then
\begin{align*}
\Px\{\tau\le t,X_t=y\}
 &=\sum_{s\le t}\sum_z\Px\{\tau=s,X_s=z,X_t=y\}\\
 &=\sum_{s\le t}\sum_z P^{t-s}(z,y)\Px\{\tau=s\}\pi(z).
\end{align*}
It follows from the stationarity of $\pi$ that
$\sum_z\pi(z)P^{t-s}(z,y)=\pi(y)$, whence for all $t\ge0$ and $y$,
\begin{equation}
\Px\{\tau\le t,X_t=y\}=\Px\{\tau\le t\}\pi(y).
\tag{6.4}\label{eq:6.4}
\end{equation}

\sourceloc{Section 6.4, physical page 95, printed page 79; separation definitions and complete Lemma 6.12}
It will be convenient to introduce a parameter $s_x(t)$, called
\emph{separation distance} and defined by
\begin{equation}
s_x(t):=\max_{y\in\X}\left[1-\frac{P^t(x,y)}{\pi(y)}\right].
\tag{6.6}\label{eq:6.6}
\end{equation}
The distance $s_x(t)$ is weakly decreasing in $t$ (see Exercise 6.4.)
We also define
\begin{equation}
s(t):=\max_{x\in\X}s_x(t).
\tag{6.7}\label{eq:6.7}
\end{equation}
The Convergence Theorem implies that $s(t)\to0$ as $t\to\infty$ for
aperiodic irreducible chains. See also Lemma 6.17.

The relationship between $s_x(t)$ and strong stationary times is

\noindent\textbf{Lemma 6.12.} \emph{If $\tau$ is a strong stationary time
for starting state $x$, then}
\begin{equation}
s_x(t)\le\Px\{\tau>t\}.
\tag{6.8}\label{eq:6.8}
\end{equation}
\begin{proof}
Fix $x\in\X$. Observe that for every $y\in\X$,
\begin{equation}
1-\frac{P^t(x,y)}{\pi(y)}
 =1-\frac{\Px\{X_t=y\}}{\pi(y)}
 \le1-\frac{\Px\{X_t=y,\tau\le t\}}{\pi(y)}.
\tag{6.9}\label{eq:6.9}
\end{equation}
By Remark 6.8, the right-hand side equals
\begin{equation}
1-\frac{\pi(y)\Px\{\tau\le t\}}{\pi(y)}=\Px\{\tau>t\}.
\tag{6.10}\label{eq:6.10}
\end{equation}
\end{proof}

\subsection*{The cited exercise and the separate author erratum}
\sourceloc{Complete original Exercise 6.4, physical page 102, printed page 86}
\noindent\textbf{Exercise 6.4.} Let $s(t)$ be defined as in \eqref{eq:6.7}.
\begin{enumerate}
\renewcommand{\labelenumi}{(\alph{enumi})}
\item Show that for each $t\ge1$ there is a stochastic matrix $Q_t$ so that
$P^t(x,\cdot)=[1-s_x(t)]\pi+s_x(t)Q_t(x,\cdot)$ and $\pi=\pi Q_t$.
\item Using the representation in (a), show that
\begin{equation}
\begin{aligned}
P^{t+u}(x,y)
 &=[1-s_x s(t)s(u)]\pi(y)\\
 &\quad+s_x(t)s(u)\sum_{z\in\X}Q_t(x,z)Q_u(z,y).
\end{aligned}
\tag{6.25}\label{eq:6.25}
\end{equation}
\item Using \eqref{eq:6.25}, establish that $s_x(t+u)\le s_x(t)s(u)$ and
deduce that $s$ is submultiplicative, i.e., $s(t+u)\le s(t)s(u)$.
\item Deduce that $s_x(t)$ is weakly decreasing in $t$.
\end{enumerate}

\sourceloc{Separate cached \emph{Errata for Markov Chains and Mixing Times, Second Edition}, physical page 1; Exercise 6.4 entry}
\begin{quote}
Exercise 6.4. The identity $\pi=\pi Q_t$ is false for $Q_t$ as defined.
Correct by replacing $s_x(t)$ by $s(t)$ in (a) and (6.25).
(From Sanchayan Bhowal.)
\end{quote}
\editorial{Exercise 6.4 remains an exercise, without an added solution.
The original exercise, including the printed first coefficient in
\eqref{eq:6.25}, is transcribed separately from the erratum. The erratum
has not been substituted into the original Proposition 6.21 proof.}
% END CONTEXT MCMT:6.21

\clearpage
\section*{6.7. Optimal Strong Stationary Times*}
\sourceloc{Physical page 101, printed page 85; original complete statement and proof}
% BEGIN PRIMARY STATEMENT MCMT:6.21
Consider an irreducible and aperiodic Markov chain $(X_t)$.

\noindent\textbf{Proposition 6.21.} \emph{For every starting state $x$,
there exists a strong stationary time $\tau$ such that, for all $t\ge0$,}
\begin{equation}
s_x(t)=\Px\{\tau>t\}.
\tag{6.20}\label{eq:6.20}
\end{equation}
% END PRIMARY STATEMENT MCMT:6.21

% BEGIN PRIMARY PROOF MCMT:6.21
\begin{proof}
Fix $x\in\X$, and let $a_t:=\min_y\frac{P^t(x,y)}{\pi(y)}=1-s_x(t)$.
Note that $a_t$ is non-decreasing. (See Exercise 6.4.) If there exists a
strong stationary time $\tau$ satisfying \eqref{eq:6.20}, then
$\Px\{\tau=t\}=a_t-a_{t-1}$ and
\begin{equation}
\Px\{X_t=y,\tau=t\}=\pi(y)(a_t-a_{t-1})
\qquad\text{for all }t\ge0,\ y\in\X.
\tag{6.21}\label{eq:6.21}
\end{equation}
Likewise, by \eqref{eq:6.4}, if \eqref{eq:6.21} holds, then
\begin{equation}
\Px\{X_t=y,\tau\le t\}=\pi(y)a_t
\qquad\text{for all }t\ge0,\ y\in\X.
\tag{6.22}\label{eq:6.22}
\end{equation}
Since
\begin{equation}
\begin{aligned}
&\Px\{X_t=y,\tau=t\}\\
&\quad=\Px\{\tau=t\mid X_t=y,\tau>t-1\}\Px\{X_t=y,\tau>t-1\}\\
&\quad=\Px\{\tau=t\mid X_t=y,\tau>t-1\}
 \big(P^t(x,y)-\Px\{X_t=y,\tau\le t-1\}\big),
\end{aligned}
\tag{6.23}\label{eq:6.23}
\end{equation}
if the optimal $\tau$ exists (so that \eqref{eq:6.21} and
\eqref{eq:6.22} are satisfied), then
\begin{equation}
\Px\{\tau=t\mid X_t=y,\tau>t-1\}
 =\frac{a_t-a_{t-1}}{P^t(x,y)/\pi(y)-a_{t-1}}.
\tag{6.24}\label{eq:6.24}
\end{equation}
The quantity on the right is in $[0,1]$ since $a_t$ is non-decreasing and
$a_t\le P^t(x,y)/\pi(y)$. To construct $\tau$ which satisfies
\eqref{eq:6.24}, let $U_1,U_2,\ldots$ be i.i.d. Uniform$[0,1]$ random
variables independent of $(X_t)$, and define
\[
\tau=\min\left\{t\ge1:\ U_t\le
\frac{a_t-a_{t-1}}{P^t(x,X_t)/\pi(X_t)-a_{t-1}}\right\}.
\]
Clearly \eqref{eq:6.24} holds. We now show the equality
\eqref{eq:6.21} holds by induction. The case $t=1$ is immediate.
For the induction step, assume that \eqref{eq:6.21} holds with $s$
replacing $t$ for all $s<t$. The proof of \eqref{eq:6.4} shows that
\[
\Px\{X_t=y,\tau\le t-1\}=\pi(y)a_{t-1}.
\]
Equation \eqref{eq:6.23} yields the induction step and then proves
\eqref{eq:6.21}. Summing \eqref{eq:6.22} over $y$ shows that
$\Px\{\tau\le t\}=a_t$ for all $t$. Consequently, \eqref{eq:6.20}
holds. In particular, $\Px\{\tau<\infty\}=1$, since $a_t\to1$ for an
aperiodic and irreducible chain. The strong stationarity of $\tau$ follows
from \eqref{eq:6.21}.
\end{proof}
% END PRIMARY PROOF MCMT:6.21
\end{document}
